\documentclass[11pt]{article}

\usepackage[utf8]{inputenc}
\usepackage[T1]{fontenc}
\usepackage[a4paper,margin=1.35cm]{geometry}
\usepackage{graphicx}
\usepackage{xcolor}
\usepackage{enumitem}
\usepackage{hyperref}
\usepackage{url}

\IfFileExists{qrcode.sty}{
  \usepackage{qrcode}
  \newcommand{\SeminarQRCode}{\qrcode[height=2.2cm]{https://alequisgs.github.io/etale-cohomology-seminar-2026-/}}
}{
  \newcommand{\SeminarQRCode}{}
}

\definecolor{Board}{HTML}{17231D}
\definecolor{Copper}{HTML}{703918}
\definecolor{Paper}{HTML}{F6EFE1}
\definecolor{Gold}{HTML}{C99B4F}

\hypersetup{
  colorlinks=true,
  urlcolor=Copper,
  linkcolor=Copper,
  pdftitle={Étale Cohomology Seminar 2026 invitation},
  pdfauthor={Étale Cohomology Seminar 2026}
}

\setlength{\parindent}{0pt}
\setlength{\parskip}{0.45em}
\setlist[itemize]{leftmargin=1.2em,itemsep=0.15em,topsep=0.2em}
\pagestyle{empty}

\begin{document}

\colorbox{Board}{%
  \begin{minipage}{\dimexpr\linewidth-2\fboxsep\relax}
    \color{white}
    \vspace{0.2cm}
    {\LARGE\bfseries Étale Cohomology Seminar 2026}\\[0.15cm]
    {\large Five weekly meetings · Online}
    \vspace{0.2cm}
  \end{minipage}
}

\vspace{0.45cm}

\includegraphics[width=\linewidth,height=4.6cm,keepaspectratio]{docs/assets/serre-grothendieck-etale-hero.png}

{\footnotesize\color{Copper}An imaginative illustration of the conversations and ideas surrounding the development of étale cohomology.}

We are organizing a weekly reading/discussion seminar following Daniel Litt's course
\href{https://www.daniellitt.com/tale-cohomology/}{\emph{Étale Cohomology and the Weil Conjectures}}.

\textbf{Starts Friday, August 28, 2026}\\
\textbf{17:00-18:30 Lima / 19:00-20:30 Brasília}

Videos, notes, and AI tools are used for preparation. Meetings are for mathematical
community: verification, questions, examples, peer review, proof digestion, and
shared understanding.

{\large\bfseries How it works: 3 + 2 + 1 + E + R + P}

\begin{itemize}
  \item \textbf{3} things you think you understood.
  \item \textbf{2} points of mathematical friction.
  \item \textbf{1} question worth discussing with humans.
  \item \textbf{E} one example or non-example.
  \item \textbf{R} review one contribution from another participant.
  \item \textbf{P} engage seriously with the common Problem of the Week.
\end{itemize}

Open to undergraduate, master's, PhD students, and other interested mathematicians.
You do not need to understand every lecture completely before coming.

\vspace{0.25cm}

\colorbox{Paper}{%
  \begin{minipage}{\dimexpr\linewidth-2\fboxsep\relax}
    \centering
    {\Large\bfseries Interested?}\\[0.1cm]
    \textbf{Read the full plan and register here:}\\[0.1cm]
    \href{https://alequisgs.github.io/etale-cohomology-seminar-2026-/}{\texttt{https://alequisgs.github.io/etale-cohomology-seminar-2026-/}}\\[0.15cm]
    Google Form registration is linked from the site.\\
    \SeminarQRCode
  \end{minipage}
}

\vspace{0.35cm}

{\large\bfseries Five pilot topics}

\begin{itemize}
  \item Weil conjectures and motivation.
  \item Étale morphisms and introduction to sites.
  \item Sites, sheaves, étale and fppf topologies.
  \item Morphisms of sites and fppf descent I.
  \item fppf descent II and the category of sheaves.
\end{itemize}

Meeting link sent to registered participants.

Funnel: WhatsApp / Telegram $\rightarrow$ seminar website $\rightarrow$ Google Form registration
$\rightarrow$ confirmation email $\rightarrow$ Google Meet + GitHub.

GitHub project:
\href{https://github.com/alequisGS/etale-cohomology-seminar-2026-}{\texttt{github.com/alequisGS/etale-cohomology-seminar-2026-}}

{\footnotesize Curriculum based on Daniel Litt's course. Seminar format developed independently by the group.}

\end{document}
